% \chapter{A few more additional information} \label{chap:appendix-4}

\section{Control Vectors Full Results}\label{appendix:ch-4-full-results}
We report the results of all the ablations in~\Cref{tab:main-results-llama}, ~\Cref{tab:main-results-gemma} and ~\Cref{tab:main-results-mistral}. Overall, we see that the baselines do not perform as well as the control vectors. We also see that some of the methods are either unsupported or cannot run due to out of memory errors. ICL, SFT, QLoRA/LoRA SFT, and prompt optimization perform poorly, with ICL and prompt optimization aligning with prior work showing they are insufficient for readability control \cite{ribeiro-etal-2023-generating}. Training-based methods also underperform, due to dataset differences (scientific texts vs. shorter news and fewer examples). 




\begin{table*}[!ht]
\centering
\scriptsize
\setlength{\extrarowheight}{1pt} \setlength{\tabcolsep}{3pt} 
\begin{subtable}[t]{0.45\textwidth}
\centering
\begin{tabular}{lll|lll}
\textbf{Method} & \textbf{CC} & \textbf{Red.} & \textbf{P} & \textbf{KT} & \textbf{BertS} \\ \hline
ICL & - & - & -0.025 & -0.017 & 0.847 \\
DSPy & - & - & 0.000 & 0.000 & 0.000 \\
QLoRa & - & - & -0.115 & -0.164 & 0.847 \\
SFT & - & - & 0.033 & 0.035 & 0.780 \\ \cdashline{1-6}[1pt/3pt]
\multirow{4}{*}{\begin{tabular}[c]{@{}l@{}}Control \\ Vectors\end{tabular}} & Diff & Avg & 0.302 & 0.205 & 0.837 \\
 & Center & Avg & 0.302 & 0.205 & 0.803 \\
 & Diff & PCA & 0.101 & 0.075 & 0.836 \\
 & Center & PCA & 0.101 & 0.075 & 0.835 \\
\end{tabular}
\caption{eLife}
\label{tab:main-results-llama-elife}
\end{subtable}
\hfill
\begin{subtable}[t]{0.45\textwidth}
\centering
\begin{tabular}{lll|lll}
\textbf{Method} & \textbf{CC} & \textbf{Red.} & \textbf{P} & \textbf{KT} & \textbf{BertS} \\ \hline
ICL & - & - & 0.00024 & -0.003 & 0.848 \\
DSPy & - & - & 0.000 & 0.000 & 0.000 \\
LoRa & - & - & -0.084 & -0.126 & 0.819 \\
SFT & - & - & 0.007 & 0.004 & 0.719 \\ \cdashline{1-6}[1pt/3pt]
\multirow{4}{*}{\begin{tabular}[c]{@{}l@{}}Control \\ Vectors\end{tabular}} & Diff & Avg & 0.416 & 0.286 & 0.805 \\
 & Center & Avg & 0.417 & 0.287 & 0.805 \\
 & Diff & PCA & 0.365 & 0.249 & 0.838 \\
 & Center & PCA & 0.144 & 0.107 & 0.837 \\
\end{tabular}
\caption{PLOS}
\label{tab:main-results-llama-plos}
\end{subtable}

\vspace{1em}

\begin{subtable}[t]{0.45\textwidth}
\centering
\begin{tabular}{lll|lll}
\textbf{Method} & \textbf{CC} & \textbf{Red.} & \textbf{P} & \textbf{KT} & \textbf{BertS} \\ \hline
ICL & - & - & 0.055 & 0.0194 & 0.853 \\
DSPy & - & - & 0.000 & 0.000 & 0.000 \\
LoRa & - & - & -0.034 & -0.036 & 0.846 \\
SFT & - & - & 0.0078 & 0.0095 & 0.764 \\ \cdashline{1-6}[1pt/3pt]
\multirow{4}{*}{\begin{tabular}[c]{@{}l@{}}Control \\ Vectors\end{tabular}} & Diff & Avg & 0.407 & 0.298 & 0.840 \\
 & Center & Avg & 0.408 & 0.298 & 0.810 \\
 & Diff & PCA & 0.062 & 0.043 & 0.839 \\
 & Center & PCA & 0.168 & 0.121 & 0.838 \\
\end{tabular}
\caption{Eureka}
\label{tab:main-results-llama-eureka}
\end{subtable}
\hfill
\begin{subtable}[t]{0.45\textwidth}
\centering
\begin{tabular}{lll|lll}
\textbf{Method} & \textbf{CC} & \textbf{Red.} & \textbf{P} & \textbf{KT} & \textbf{BertS} \\ \hline
ICL & - & - & 0.0302 & -0.004 & 0.840 \\
DSPy & - & - & 0.000 & 0.000 & 0.000 \\
LoRa & - & - & -0.048 & -0.053 & 0.480 \\
SFT & - & - & -0.0062 & -0.0059 & 0.44 \\ \cdashline{1-6}[1pt/3pt]
\multirow{4}{*}{\begin{tabular}[c]{@{}l@{}}Control \\ Vectors\end{tabular}} & Diff & Avg & 0.075 & 0.000 & 0.827 \\
 & Center & Avg & 0.268 & 0.216 & 0.801 \\
 & Diff & PCA & 0.395 & 0.315 & 0.826 \\
 & Center & PCA & 0.081 & 0.074 & 0.826 \\
\end{tabular}
\caption{SciNews}
\label{tab:main-results-llama-scinews}
\end{subtable}
\caption[Full set of results for all the Control Vector settings we tested for \texttt{meta-llama/Llama-3.1-8B-Instruct}.]{Full set of results for all the Control Vector settings we tested for \texttt{meta-llama/Llama-3.1-8B-Instruct}. We report the Pearson (P) and Kendall-Tau (KT) correlation of the specified level of readability with the Coleman-Liau readability index, along with BertScore F1 (BertS). We experiment with In-Context Learning (ICL), Supervised Finetuning (SFT), QLoRa/LoRa SFT, and DSPy Prompt Optimization (DSPy) as our baselines. For the Control Vectors (CV), we experiment with 2 Contrastive Combination (CC) methods (Diff and Center) and 2 Reduction (Red.) methods (Avg. and PCA).}
\label{tab:main-results-llama}
\end{table*}


\begin{table*}[!ht]
\centering
\scriptsize
\setlength{\extrarowheight}{1pt} \setlength{\tabcolsep}{3pt} 
\begin{subtable}[t]{0.45\textwidth}
\centering
\begin{tabular}{lll|lll}
\textbf{Method} & \textbf{CC} & \textbf{Red.} & \textbf{P} & \textbf{KT} & \textbf{BertS} \\ \hline
DSPy & - & - & \multicolumn{3}{c}{CONTEXT ERROR} \\
LoRa & - & - & \multicolumn{3}{c}{OOM} \\
SFT & - & - & \multicolumn{3}{c}{OOM} \\ \cdashline{1-6}[1pt/3pt]
\multirow{4}{*}{\begin{tabular}[c]{@{}l@{}}Control \\ Vectors\end{tabular}} & Diff & Avg & 0.000 & 0.000 & 0.822 \\
 & Center & Avg & 0.132 & 0.206 & 0.808 \\
 & Diff & PCA & 0.315 & 0.223 & 0.821 \\
 & Center & PCA & 0.280 & 0.207 & 0.823 \\
\end{tabular}
\caption{eLife}
\label{tab:main-results-gemma-elife}
\end{subtable}
\hfill
\begin{subtable}[t]{0.48\textwidth}
\centering
\begin{tabular}{lll|lll}
\textbf{Method} & \textbf{CC} & \textbf{Red.} & \textbf{P} & \textbf{KT} & \textbf{BertS} \\ \hline
DSPy & - & - & \multicolumn{3}{c}{CONTEXT ERROR} \\
LoRa & - & - & 0.000 & 0.000 & 0.000 \\
SFT & - & - & \multicolumn{3}{c}{OOM} \\ \cdashline{1-6}[1pt/3pt]
\multirow{4}{*}{\begin{tabular}[c]{@{}l@{}}Control \\ Vectors\end{tabular}} & Diff & Avg & 0.000 & 0.000 & 0.832 \\
 & Center & Avg & 0.085 & 0.112 & 0.815 \\
 & Diff & PCA & 0.213 & 0.189 & 0.830 \\
 & Center & PCA & 0.100 & 0.120 & 0.831 \\
\end{tabular}
\caption{PLOS}
\label{tab:main-results-gemma-plos}
\end{subtable}

\vspace{1em}

\begin{subtable}[t]{0.45\textwidth}
\centering
\begin{tabular}{lll|lll}
\textbf{Method} & \textbf{CC} & \textbf{Red.} & \textbf{P} & \textbf{KT} & \textbf{BertS} \\ \hline
DSPy & - & - & \multicolumn{3}{c}{CONTEXT ERROR} \\
LoRa & - & - & -0.002 & 0.009 & 0.813 \\
SFT & - & - & 0.003 & -0.011 & 0.798 \\ \cdashline{1-6}[1pt/3pt]
\multirow{4}{*}{\begin{tabular}[c]{@{}l@{}}Control \\ Vectors\end{tabular}} & Diff & Avg & 0.000 & 0.000 & 0.906 \\
 & Center & Avg & 0.117 & 0.147 & 0.870 \\
 & Diff & PCA & 0.137 & 0.112 & 0.891 \\
 & Center & PCA & 0.121 & 0.098 & 0.899 \\
\end{tabular}
\caption{Eureka}
\label{tab:main-results-gemma-eureka}
\end{subtable}
\hfill
\begin{subtable}[t]{0.45\textwidth}
\centering
\begin{tabular}{lll|lll}
\textbf{Method} & \textbf{CC} & \textbf{Red.} & \textbf{P} & \textbf{KT} & \textbf{BertS} \\ \hline
DSPy & - & - & \multicolumn{3}{c}{CONTEXT ERROR} \\
LoRa & - & - & -0.034 & -0.040 & 0.483 \\
SFT & - & - & \multicolumn{3}{c}{OOM} \\ \cdashline{1-6}[1pt/3pt]
\multirow{4}{*}{\begin{tabular}[c]{@{}l@{}}Control \\ Vectors\end{tabular}} & Diff & Avg & 0.000 & 0.000 & 0.824 \\
 & Center & Avg & 0.141 & 0.177 & 0.809 \\
 & Diff & PCA & 0.116 & 0.110 & 0.822 \\
 & Center & PCA & 0.120 & 0.114 & 0.824 \\
\end{tabular}
\caption{SciNews}
\label{tab:main-results-gemma-scinews}
\end{subtable}
\caption[Full set of results for all the Control Vector settings we tested for \texttt{google/gemma-7b-it}.]{Full set of results for all the Control Vector settings we tested for \texttt{google/gemma-7b-it}. We report the Pearson (P) and Kendall-Tau (KT) correlation of the specified level of readability with the Coleman-Liau readability index, along with BertScore F1 (BertS). We experiment with In-Context Learning (ICL), Supervised Finetuning (SFT), QLoRa/LoRa SFT, and DSPy Prompt Optimization (DSPy) as our baselines. For the Control Vectors (CV), we experiment with 2 Contrastive Combination (CC) methods (Diff and Center) and 2 Reduction (Red.) methods (Avg. and PCA). OOM means the baseline could not be run due to an out of memory error on our setup. UNSUPPORTED means that out of the box the baseline does not support \texttt{google/gemma-7b-it}.}
\label{tab:main-results-gemma}
\end{table*}






\begin{table*}[!ht]
\centering
\scriptsize
\setlength{\extrarowheight}{1pt} \setlength{\tabcolsep}{3pt} 
\begin{subtable}[t]{0.45\textwidth}
\centering
\begin{tabular}{lll|lll}
\textbf{Method} & \textbf{CC} & \textbf{Red.} & \textbf{P} & \textbf{KT} & \textbf{BertS} \\ \hline
DSPy & - & - & \multicolumn{3}{c}{CONTEXT ERROR} \\
LoRa & - & - & 0.001 & 0.004 & 0.850 \\
SFT & - & - & 0.001 & 0.000 & 0.724 \\ \cdashline{1-6}[1pt/3pt]
\multirow{4}{*}{\begin{tabular}[c]{@{}l@{}}Control \\ Vectors\end{tabular}} & Diff & Avg & 0.000 & 0.000 & 0.822 \\
 & Center & Avg & 0.132 & 0.206 & 0.808 \\
 & Diff & PCA & 0.315 & 0.223 & 0.821 \\
 & Center & PCA & 0.280 & 0.207 & 0.823 \\
\end{tabular}
\caption{eLife}
\label{tab:main-results-mistral-elife}
\end{subtable}
\hfill
\begin{subtable}[t]{0.45\textwidth}
\centering
\begin{tabular}{lll|lll}
\textbf{Method} & \textbf{CC} & \textbf{Red.} & \textbf{P} & \textbf{KT} & \textbf{BertS} \\ \hline
DSPy & - & - & \multicolumn{3}{c}{CONTEXT ERROR} \\
LoRa & - & - & 0.000 & 0.001 & 0.857 \\
SFT & - & - & 0.003 & 0.005 & 0.424 \\ \cdashline{1-6}[1pt/3pt]
\multirow{4}{*}{\begin{tabular}[c]{@{}l@{}}Control \\ Vectors\end{tabular}} & Diff & Avg & 0.000 & 0.000 & 0.832 \\
 & Center & Avg & 0.085 & 0.112 & 0.815 \\
 & Diff & PCA & 0.213 & 0.189 & 0.830 \\
 & Center & PCA & 0.100 & 0.120 & 0.831 \\
\end{tabular}
\caption{PLOS}
\label{tab:main-results-mistral-plos}
\end{subtable}

\vspace{1em}

\begin{subtable}[t]{0.45\textwidth}
\centering
\begin{tabular}{lll|lll}
\textbf{Method} & \textbf{CC} & \textbf{Red.} & \textbf{P} & \textbf{KT} & \textbf{BertS} \\ \hline
DSPy & - & - & \multicolumn{3}{c}{CONTEXT ERROR} \\
LoRa & - & - & 0.027 & 0.035 & 0.766 \\
SFT & - & - & 0.000 & 0.000 & 0.696 \\ \cdashline{1-6}[1pt/3pt]
\multirow{4}{*}{\begin{tabular}[c]{@{}l@{}}Control \\ Vectors\end{tabular}} & Diff & Avg & 0.000 & 0.000 & 0.906 \\
 & Center & Avg & 0.117 & 0.147 & 0.870 \\
 & Diff & PCA & 0.137 & 0.112 & 0.891 \\
 & Center & PCA & 0.121 & 0.098 & 0.899 \\
\end{tabular}
\caption{Eureka}
\label{tab:main-results-mistral-eureka}
\end{subtable}
\hfill
\begin{subtable}[t]{0.45\textwidth}
\centering
\begin{tabular}{lll|lll}
\textbf{Method} & \textbf{CC} & \textbf{Red.} & \textbf{P} & \textbf{KT} & \textbf{BertS} \\ \hline
DSPy & - & - & \multicolumn{3}{c}{CONTEXT ERROR} \\
LoRa & - & - & 0.000 & 0.000 & 0.104 \\
SFT & - & - & 0.000 & 0.000 & 0.361 \\ \cdashline{1-6}[1pt/3pt]
\multirow{4}{*}{\begin{tabular}[c]{@{}l@{}}Control \\ Vectors\end{tabular}} & Diff & Avg & 0.000 & 0.000 & 0.824 \\
 & Center & Avg & 0.141 & 0.177 & 0.809 \\
 & Diff & PCA & 0.116 & 0.110 & 0.822 \\
 & Center & PCA & 0.120 & 0.114 & 0.824 \\
\end{tabular}
\caption{SciNews}
\label{tab:main-results-mistral-scinews}
\end{subtable}
\caption[Full set of results for all the Control Vector settings we tested for \texttt{mistralai/Mistral-7B-Instruct-v0.3}.]{Full set of results for all the Control Vector settings we tested for \texttt{mistralai/Mistral-7B-Instruct-v0.3}. We report the Pearson (P) and Kendall-Tau (KT) correlation of the specified level of readability with the Coleman-Liau readability index, along with BertScore F1 (BertS). We experiment with In-Context Learning (ICL), Supervised Finetuning (SFT), QLoRa/LoRa SFT, and DSPy Prompt Optimization (DSPy) as our baselines. For the Control Vectors (CV), we experiment with 2 Contrastive Combination (CC) methods (Diff and Center) and 2 Reduction (Red.) methods (Avg. and PCA). OOM means the baseline could not be run due to an out of memory error on our setup. LENGTH ERROR means the baseline could not be run due to the context window being too short.}
\label{tab:main-results-mistral}
\end{table*}

\subsection{Baselines for Control Vectors Discussion}\label{appendix:chap-4-baseline-discussion}
The baselines were less simple to run than the control vectors leading to four major issues: Resource Usage, Out of Memory Errors, Lack of Support Out of the Box, and Small Training Set. 
\paragraph{Resource Usage}
For DSPy, it also used way more resources than control vectors while performing worse. The DSPy method we use uses in context learning which forces the prompt to include extremely long examples. This could lead to some of the sources not being able to be summarized as the total length of the prompt was more than the max context length of the model. This was the main issue with DSPy for these smaller models. 

As for SFT and QLoRa/LoRa-SFT, it requires more forward passes as finetuning is multiple epochs with the same amount of tokens. In addition, it requires expensive backward passes that require more memory. Even with this more computation, our method performs better while keeping training costs down. 
\paragraph{Out of Memory Errors}
Many of the training based methods required hefty amounts of resources and very long source documents led to out of memory errors with even QLoRa. While truncation could have prevented these errors, truncating at an arbitrary length could harm the overall plain language summary as this increases problems with reading comprehension, adds more hallucinations, and more \cite{ding2024fewertruncationsimprovelanguage}. While smarter engineering could have performed better with other training optimizations, the amount of work to prevent these errors would have dwarfed the time it takes to set up the control vectors and could have led to even slower training times with the increase in communication time between devices and loading from RAM \cite{rajbhandari2020zeromemoryoptimizationstraining}. In addition, Gemma experienced more out of memory errors with fewer parameters  due to a different less memory efficient architecture. 
\paragraph{Lack of Support Out of the Box}
We could not run DSPy prompt optimization on the Gemma model due to an incompatibility between DSPy and the Gemma system prompt. It would have required modifying DSPy for a specific set of models by adding additional adapters not available out of the box. 
\paragraph{Small Training Set}
We hypothesize that SFT performed worse in controllability than LoRA due to the number of trainable parameters relative to the size of the training dataset. Our SFT model was trained on only the 128 most readable examples, which likely provided insufficient signal to effectively update the large number of parameters involved in full fine-tuning. In contrast, LoRA introduces far fewer trainable parameters, making it better suited for small-data regimes. This is consistent with the finding in \cite{whitehouse-etal-2024-low} where they find that "in low-data scenarios, LoRA is a better alternative to full fine-tuning."

At first glance, this result may appear to contradict the findings of~\textcite{ribeiro-etal-2023-generating}. However, their experiments were conducted with a substantially larger training dataset. Additionally, their source documents consisted primarily of news articles, which are typically much shorter and less structurally complex than scientific articles, reducing the difficulty of the summarization task.